\quellseite{3}
\begin{titelseite}
\werkuntertitel{Works of Prof. G. H. Gerberding, D.D., LL.D.}

\emph{THE WAY OF SALVATION IN THE LUTHERAN CHURCH}\\
(Thirtieth Thousand)\\
The most popular English Lutheran Book. Has been translated into German, Swedish, Norwegian, and in part into Tamil, Telugu, and Japanese. Has converted scores of outsiders to the Lutheran Faith. The best seller. In its 30th 1000 in English. 267 pages.

Price, \$1.00

\emph{NEW TESTAMENT CONVERSIONS}\\
(Seventh Edition)\\
Plain. Practical. Popular sermons, pervaded by a spirit of deep piety. Meets and answers “Salvation by Character” and emotional revivalism by an earnest appeal to apostolic teaching and practice. In its seventh edition.

Price, \$1.00

\emph{LIFE AND LETTERS OF W. A. PASSAVANT, D. D.}\\
(Second Edition)\\
The life story, largely autobiographical, of the most remarkable man in the American Lutheran Church. It is, indeed, doubtful whether any other church has produced his like. This life story is a revelation and an inspiration. Every Lutheran who loves his Church and would learn to understand her better and love her more ought to read it.

Over 600 pages. Price, \$2.00

\emph{THE LUTHERAN PASTOR}\\
(Seventh Edition)\\
A Text-book on Pastoral Theology. The only English work on the subject on the market.

Highly commended by \emph{The Lutheran}, \emph{The Lutheran World}, \emph{Observer}, \emph{The Lutheran Standard}, \emph{The Lutheran Witness}, Dr. Stellhorn in \emph{Theologische Zeitblaetter}, President Richter in \emph{Iowa Kirchen-Blatt}, Prof. Geo. H. Schodde and ministers from every part of our Lutheran Church.

462 pages. Price, \$2.00 net

\emph{THE LUTHERAN CATECHIST}\\
(Fourth Edition)\\
A Companion Book to “The Lutheran Pastor”

The Publisher’s Announcement says: “A book for the pastor, whose difficulties, cares, needs and joys Dr. Gerberding so well understands. It will do good to every student of theology and be a great blessing to the mature minister. By its perusal the usefulness of many catechists may be doubled.”

Few new books have been received with such general favor.

Over 300 pages. Price, \$1.50 net

In Preparation\\
\emph{THE LUTHERAN CHURCH IN THE COUNTRY}
\end{titelseite}

\quellseite{6}
\begin{titelseite}
\werktitel{THE LUTHERAN CATECHIST}

\werkuntertitel{A Companion Book to “The Lutheran Pastor.”}

by

G. H. Gerberding, D.D., LL.D.,

Professor of Practical Theology in the Theological Seminary of the Evangelical Lutheran Church, Chicago,

author of

“The Way of Salvation in the Lutheran Church,”\\
“New Testament Conversions,” “The Lutheran Pastor,”\\
“Life and Letters of Passavant,” etc.

Fourth Revised Edition.

The Lutheran Publication Society,

Philadelphia.
\end{titelseite}

\quellseite{7}
\begin{titelseite}
[Copyright, 1910, by G. H. Gerberding.]

\werktitel{Dedication.}

In order that the Church of his love, in the future, may hold, mould and make more Christ-like the multitude of promising children that God gives her, the author dedicates this book

\werkuntertitel{TO THE CAUSE OF CHRISTIAN EDUCATION}

In which the Church of the Reformation has always been the leader in other lands, and in which she must lead in this land also if she would take and hold that place of leadership to which she is entitled by her history, her genius, her spirit, her type of piety, and her evangelical doctrine, all of which are so beautifully fitted for the closet, the family, the Sunday school, the catechetical class and the pulpit.
\end{titelseite}

\quellseite{8}
\chapter[Preface to the Fourth Revised Edition]{PREFACE TO THE FOURTH REVISED EDITION}

The writer is grateful to the dear Lord and to an appreciative Church for the kind reception accorded to the Lutheran Catechist. It has had a wide circulation among Lutheran ministers in all the divisions of the Church. It is in use in nearly all the Theological Seminaries in our Church, either as a text-book or as a companion book of reference. The commendations have been most hearty and grateful. New editions are being called for in quick succession.

In sending forth this fourth revised edition, the writer would acknowledge one rather severe review by the Rev. Prof. Reu, of Wartburg Theological Seminary, Dubuque, Iowa. In this book the writer acknowledges the superior scholarship of Prof. Reu, and pays deserved tribute to his learning. And the writer here frankly and gladly acknowledges that Dr. Reu knows far more about the history of Christian instruction than he.

A few historic statements have been changed. The slips, however, had not been serious.

Dr. Reu wrongly accuses the writer of knowing nothing of the need and value of Bible History in
\quellseite{9}
Catechisation. The writer might make the retort courteous that the learned doctor knows nothing of our excellent graded Sunday school series, which has the best eight years’ course ever worked out in Bible Story, History, Readings, Facts and Scenes, Biography, Teachings, Literature and Outlines. This series is mentioned several times in the Catechist in its relation to catechetical instruction.

We did not expect Dr. Reu to agree with us on the Public Examination, or on the question and answer form of teaching. These are matters of expediency on which good Lutherans may differ.

In criticising Part V, Dr. Reu forgets that these are “Hints” for the teacher, and not expositions of the five parts.

We send out this new edition with the same hope and prayer with which we sent out the first.

\begin{unterschrift}
Chicago Seminary, October, 1915.
\end{unterschrift}

\quellseite{10}
\chapter{PREFACE}

The Lutheran minister has three special activities.\randnotiz{The minister’s three special activities.} He is, first of all, a preacher of the Word, the whole Word, its Law and its Gospel. He is always and everywhere a \emph{seelsorger}, a pastor, an under-shepherd, seeking the lost, feeding and caring for those in the fold, knowing them by name, carrying on his heart their woes, their wants and their welfare, looking after them individually and collectively. If he would make full proof of his ministry, he must love the Shepherd and Bishop of his own soul, and must show his love by feeding His lambs, by tending His young sheep. He must be a good Catechist.

Each one of these ministerial functions requires special gifts and special preparation. Our theological seminaries deal with young men who are supposed to have the gifts, some in smaller, others in larger measure. The seminaries try to develop and perfect these gifts. They try to give the preparation needed for the several pastoral functions. The aim of the whole course of study is to fit out ministers approved of God, workmen that need not be ashamed. Much attention is given to homiletical study and drill so that each graduate may go out
\quellseite{11}
and be an attractive and a persuasive preacher of Christ.

\randnotiz{Pastoral theological literature.} Pastoral theology is given more or less attention. The purpose is to enable each student to become an active and an efficient \emph{seelsorger}. In this department there is still room for much improvement, and we all have our \emph{pia desideria}. There is need of clinical experience, demonstration and practice, under the eyes and at the side of an ideal pastor.

The man who would learn to preach, to preach well, to preach with power and prevail, has a library of books, a long line of inspiring examples in biography, and in the pulpits of the day to instruct, inspire and aid him.

He who would become an efficient, an uplifting, a health- and hope- and life-bearing pastor, a real soul curate, also has a literature and a list of living examples to teach and direct him. The English Lutheran minister can have these aids in his own tongue.

But when we come to the function and activity of the Catechist, especially the English Lutheran Catechist, the case is different. There is no English textbook for our seminaries. The professor must resort to his own lectures on Catechetics. This is a misfortune for the students. No matter how good the lectures may be, they can never take the place of a
\quellseite{12}
good text-book in the hands of student and teacher.\footnote{Along with the use of a good text-book the student should make original research and do Seminar work.} \randnotiz{Lack of catechetical literature.} The lecture system is fortunately going out of date. It may be all right as a theory. It is not a theory, but a condition, that confronts us. The student can never learn as much by listening to a course of lectures as from a text-book in hand, and afterwards at hand for reference and review. If listening were enough, what intelligent Christians, what theologians the hearers of good sermons would become! But do they?

The lecturers on Catechetics often spend a large portion of their time on the history of catechising and the history of Catechisms. This is good, but this alone will not teach and help the student rightly to catechise. Much of the remaining time is used in rehearsing the theories of German writers on the subject. This, also, is interesting and profitable as information. But German theories do not fit American conditions. And so the complaint has been long and loud that young men enter the ministry and do not know how to catechise. They feel it; they lament it. They ask for a remedy.

We cannot direct the English-speaking pastor to a literature. The books on Catechetics are in the German and Scandinavian tongues. Many of them are richly instructive and full of valuable suggestions.
\quellseite{13}
But they do not help him who knows the English language alone. \randnotiz{Books on catechetics.} There is an unending list of explanations of Luther’s Catechism, a goodly number of which have been prepared especially for the use and help of the pastor. Many of these are in English. As we are about to go to press we welcome “The Catechist’s Handbook,” by our friend, the Rev. J. W. Horine. Every English Lutheran pastor ought to have and use it, and will get much help from it in the preparation of his lessons.

But the English Catechist may still be in want. He needs a plain and practical discussion on the science and art of catechising. There is no English book now on the market that covers this ground. We believe that here is a felt want. To supply this want we have written this book. It is not an explanation of Luther’s Catechism. It is intended to help the Catechist to so use the explanations, that his work may be more satisfactory, more fruitful and more filled with delight, both for himself and for his Catechumens.

\randnotiz{This book an outgrowth.}The book is the outgrowth of nearly twenty years in the active ministry, and of fifteen years spent in teaching prospective Catechists in practical theology. In our course in Catechetics we have spent much time on practical drill work, and given the widest liberty for questions and class discussions. We believe this to be a feature of all good
\quellseite{14}
teaching for advanced students. The spare hours of three summer vacations have been employed in the writing of the book.

We do not expect every reader to agree with all the positions taken in the book. In every book that deals with human conditions, relations and duties the personal element will creep in. It can be kept out only by making the work so objective, abstract and theoretical that it loses that human element which must be the life and appeal of every book that would have a wide and varied reading.

As will be apparent from its pages we have consulted and used what we consider the best German works on the subject of Catechetics. We are indebted to the Rev. J. R. E. Hunt, of Austin, Chicago, who has made a specialty of child-psychology, for suggestions on the chapter that deals with child-study, as also for the preparation of the index.

Our hope and prayer are expressed in the dedication, found on a previous page.

\begin{unterschrift}
Chicago Seminary,\\
Lent, 1910.
\end{unterschrift}
